\section{Theoretical Consequences of Bias}
\label{sec:theory}

Section~\ref{sec:formal_model} introduced a structured decomposition of token
representations:
\[
x_i = \mu + s_i^{\text{spike}} + s_i^{\text{tail}},
\]
where $\mu$ is a shared bias component, $s_i^{\text{spike}}$ lies in a
low-dimensional subspace, and $s_i^{\text{tail}}$ occupies a high-dimensional
complement.
Section~\ref{sec:empirical} established empirically that activation bias
dominates extreme entries and significantly inflates optimizer statistics.
We now formalize these observations.

Throughout this section, we assume:
\begin{itemize}
\item $\mu$ is dense and fixed across tokens,
\item $s_i$ are zero-mean random vectors with finite second moment,
\item $s_i$ have sub-Gaussian tails (for concentration results).
\end{itemize}

The assumptions reflect the empirical statistics observed in Section~2.

\subsection{Bias Dominates Extreme Activation Entries}
\label{subsec:theory_outliers}

We begin by formalizing why bias dominates the largest-magnitude activation
entries.

\begin{theorem}[Bias dominance of extreme entries]
\label{thm:bias_outlier}
Let $x_i = \mu + s_i$ with $\mathbb{E}[s_i]=0$ and
$\|s_{ij}\|$ sub-Gaussian with variance proxy $\sigma^2$.
If $\max_j |\mu_j| \gg \sigma$,
then for sufficiently small $\varepsilon>0$, among the top
$\varepsilon$-fraction of entries (by $|x_{ij}|$),
\[
\frac{|M_{ij}|^2}{|x_{ij}|^2} \to 1
\quad \text{with high probability as } N\to\infty.
\]
\end{theorem}

\paragraph{Interpretation.}
Extreme activation entries are primarily caused by the coherent bias component.
Random fluctuations from $s_i$ produce heavy tails only with exponentially small
probability, whereas the bias contributes deterministically to every token.

\paragraph{Connection to Section~2.2.}
The empirical result that mean contribution exceeds $85\%$ of squared magnitude
for top-$0.1\%$ entries (Section~2.2) is a finite-sample manifestation of
Theorem~\ref{thm:bias_outlier}.

\subsection{Bias as a Structured Component of the Spectral Spike}
\label{subsec:theory_spike}

The activation matrix can be written as
\[
X = M + S,
\qquad
M = \mathbf{1}\mu^\top.
\]

\begin{theorem}[Bias contributes a rank-one spike]
\label{thm:bias_spike}
The matrix $M$ has rank one with singular value
\[
\sigma_1(M) = \sqrt{N}\|\mu\|_2.
\]
Moreover,
\[
\sigma_1(X) \ge \sigma_1(M) - \|S\|_2.
\]
If $\sqrt{N}\|\mu\|_2$ dominates $\|S\|_2$,
then the leading singular direction of $X$ aligns with $\mu$.
\end{theorem}

\paragraph{Interpretation.}
Bias produces a coherent low-rank component inside the spike subspace.
Mean removal eliminates this rank-one contribution exactly, but does not
eliminate structured variance in $S_{\text{spike}}$.
This explains why Section~2.3 observed partial but not complete spike reduction
after centering.

\subsection{Bias-Induced Gradient Structure}
\label{subsec:theory_grad}

Consider a linear layer with gradient
\[
G = \mathbb{E}[\delta x^\top].
\]
Substituting $x = \mu + s$ yields the exact identity
\begin{equation}
\label{eq:G_decomp_theory}
G =
\mathbb{E}[\delta]\mu^\top
+
\mathbb{E}[\delta s^\top].
\end{equation}

\begin{theorem}[Bias-driven gradient component]
\label{thm:grad_bias}
Even if $\mathbb{E}[\delta]=0$, the second-moment statistics of $G$
contain bias-dependent inflation terms:
\[
v = \mathbb{E}[g \odot g]
= \frac{1}{N}(\Delta^{\odot 2})^\top (X^{\odot 2}).
\]
If $X = M + S$, then
\[
X^{\odot 2}
=
M^{\odot 2}
+
2(M \odot S)
+
S^{\odot 2}.
\]
The first term contributes deterministic inflation proportional to $\mu_j^2$.
\end{theorem}

\paragraph{Interpretation.}
Bias increases extreme entries of $v$, shrinking effective learning rates in
Adam-type optimizers.
Mean removal reduces this inflation multiplicatively, consistent with the
empirical reductions observed in Section~2.4B.

\subsection{Dynamic-Range Inflation Under Bias}
\label{subsec:theory_quant}

Blockwise quantization uses scales proportional to
\[
\alpha = \max_{j \in \text{block}} |x_{ij}|.
\]

\begin{theorem}[Bias-driven dynamic range inflation]
\label{thm:quant_bias}
Under the model $x_i=\mu+s_i$,
if $|\mu_j|$ dominates typical fluctuations of $s_{ij}$,
then for blocks containing dimension $j$,
\[
\alpha \approx |\mu_j|
\]
with high probability.
\end{theorem}

\paragraph{Interpretation.}
Bias inflates quantization scale and compresses semantic resolution.
Mean removal restores dynamic range symmetry, explaining the empirical
improvements observed in Section~2.4C.

\paragraph{Summary of theoretical implications.}

The preceding results establish three mechanisms:

\begin{enumerate}
\item Bias dominates extreme activation entries.
\item Bias contributes a coherent low-rank spectral component.
\item Bias inflates second-moment and quantization statistics.
\end{enumerate}

These mechanisms collectively explain why mean removal is an effective
and computationally trivial intervention.
Detailed proofs of Theorems~\ref{thm:bias_outlier}--\ref{thm:quant_bias}
are provided in Appendix~A.

